% GENERATED FILE. Edit canonical lessons, then run npm run generate:books.
% canonical-source-hash: fnv1a64:aa96d79ce5ed9bc3
% canonical-lessons: LA-C90-equus, LA-C90-bos, LA-C90-ovis, LA-C90-capra, LA-C90-porcus

\chapter{Horse, Ox, Sheep, Goat, Pig}
\label{ch:la-horse-ox-sheep-goat-pig}
\hlchaptermodality{90}

\begin{chapteropening}
\textbf{By the end of this chapter:} \emph{I can name a horse, an ox, a sheep, a goat and a pig.}

Complete the last lesson of chapter 90: I can name a horse, an ox, a sheep, a goat and a pig.
\end{chapteropening}

\section[equus, equī]{equus, equī — a horse}

\label{lesson:LA-C90-equus}



\begin{quote}
Before the new one: say the Latin for a plain, then the Latin for a shore.
\end{quote}

\subsection*{You'll want to know: equus, equī}
\textbf{equus, equī} — \textquotedblleft{}a horse\textquotedblright{}. English \emph{equestrian}.

\begin{cousinweb}
\textbf{equus} is inherited. Old Latin wrote \emph{equos}, and it is usually traced to Proto-Indo-European \emph{*h1ekwos}, \textquotedblleft{}horse\textquotedblright{} (EtymDB 2.1); Greek \emph{hippos}, \textquotedblleft{}horse\textquotedblright{}, is a cousin (Lewis \& Short).
\end{cousinweb}

\begin{grammarlens}[title={What you've built}]
One more word for this chapter.
\end{grammarlens}

\subsection*{Your turn}
\begin{itemize}
\raggedright
  \item \emph{Say it:} \emph{equus}
  \item \emph{Say it:} \emph{equus}, once more
  \item \emph{Say it:} say \emph{equus}
  \item \emph{Recall:} say the Latin for a plain, then the Latin for a shore, then say \emph{equus} again
\end{itemize}

\begin{tcolorbox}[breakable,colback=teal!4,colframe=teal!35!black,title={Before you move on}]
What does equus, equī mean? (\textquotedblleft{}A horse\textquotedblright{}.) Say it once more.
\end{tcolorbox}
\section[bōs, bovis]{bōs, bovis — an ox, a cow}

\label{lesson:LA-C90-bos}



\begin{quote}
Before the new one: say the Latin for a shore, then the Latin for a horse.
\end{quote}

\subsection*{You'll want to know: bōs, bovis}
\textbf{bōs, bovis} — \textquotedblleft{}an ox, a cow\textquotedblright{}. English \emph{bovine}.

\begin{cousinweb}
\textbf{bōs} is usually traced to Proto-Indo-European \emph{*gwows}, \textquotedblleft{}cow\textquotedblright{} (EtymDB 2.1), but its b suggests it reached Latin from a neighbouring Italic language such as Oscan or Umbrian: Latin's own sound changes turned that gw into v. English \emph{cow} is a cousin.
\end{cousinweb}

\begin{grammarlens}[title={What you've built}]
One more word for this chapter.
\end{grammarlens}

\subsection*{Your turn}
\begin{itemize}
\raggedright
  \item \emph{Say it:} \emph{bōs}
  \item \emph{Say it:} \emph{bōs}, once more
  \item \emph{Say it:} say \emph{bōs}
  \item \emph{Recall:} say the Latin for a shore, then the Latin for a horse, then say \emph{bōs} again
\end{itemize}

\begin{tcolorbox}[breakable,colback=teal!4,colframe=teal!35!black,title={Before you move on}]
What does bōs, bovis mean? (\textquotedblleft{}An ox, a cow\textquotedblright{}.) Say it once more.
\end{tcolorbox}
\section[ovis, ovis]{ovis, ovis — a sheep}

\label{lesson:LA-C90-ovis}



\begin{quote}
Before the new one: say the Latin for a horse, then the Latin for an ox.
\end{quote}

\subsection*{You'll want to know: ovis, ovis}
\textbf{ovis, ovis} — \textquotedblleft{}a sheep\textquotedblright{}.

\begin{cousinweb}
\textbf{ovis} is inherited, usually traced to Proto-Indo-European \emph{*h2owis}, \textquotedblleft{}sheep\textquotedblright{} (EtymDB 2.1). English \emph{ewe} is a cousin.
\end{cousinweb}

\begin{grammarlens}[title={What you've built}]
One more word for this chapter.
\end{grammarlens}

\subsection*{Your turn}
\begin{itemize}
\raggedright
  \item \emph{Say it:} \emph{ovis}
  \item \emph{Say it:} \emph{ovis}, once more
  \item \emph{Say it:} say \emph{ovis}
  \item \emph{Recall:} say the Latin for a horse, then the Latin for an ox, then say \emph{ovis} again
\end{itemize}

\begin{tcolorbox}[breakable,colback=teal!4,colframe=teal!35!black,title={Before you move on}]
What does ovis, ovis mean? (\textquotedblleft{}A sheep\textquotedblright{}.) Say it once more.
\end{tcolorbox}
\section[capra, caprae]{capra, caprae — a goat}

\label{lesson:LA-C90-capra}



\begin{quote}
Before the new one: say the Latin for an ox, then the Latin for a sheep.
\end{quote}

\subsection*{You'll want to know: capra, caprae}
\textbf{capra, caprae} — \textquotedblleft{}a goat\textquotedblright{}.

\begin{cousinweb}
\textbf{capra} is the feminine of \emph{caper}, \textquotedblleft{}a he-goat\textquotedblright{} (Lewis \& Short), which is usually traced to Proto-Indo-European \emph{*kapros}, \textquotedblleft{}buck, he-goat\textquotedblright{} (EtymDB 2.1).
\end{cousinweb}

\begin{grammarlens}[title={What you've built}]
One more word for this chapter.
\end{grammarlens}

\subsection*{Your turn}
\begin{itemize}
\raggedright
  \item \emph{Say it:} \emph{capra}
  \item \emph{Say it:} \emph{capra}, once more
  \item \emph{Say it:} say \emph{capra}
  \item \emph{Recall:} say the Latin for an ox, then the Latin for a sheep, then say \emph{capra} again
\end{itemize}

\begin{tcolorbox}[breakable,colback=teal!4,colframe=teal!35!black,title={Before you move on}]
What does capra, caprae mean? (\textquotedblleft{}A goat\textquotedblright{}.) Say it once more.
\end{tcolorbox}
\section[porcus, porcī]{porcus, porcī — a pig}

\label{lesson:LA-C90-porcus}



\begin{quote}
Before the new one: say the Latin for a sheep, then the Latin for a goat.
\end{quote}

\subsection*{You'll want to know: porcus, porcī}
\textbf{porcus, porcī} — \textquotedblleft{}a pig\textquotedblright{}. English \emph{pork}.

\begin{cousinweb}
\textbf{porcus} is inherited, usually traced to Proto-Indo-European \emph{*porkos}, \textquotedblleft{}young pig\textquotedblright{} (EtymDB 2.1). English \emph{farrow}, \textquotedblleft{}a litter of pigs\textquotedblright{}, is a cousin.
\end{cousinweb}

\begin{grammarlens}[title={What you've built}]
One more word for this chapter.
\end{grammarlens}

\subsection*{Your turn}
\begin{itemize}
\raggedright
  \item \emph{Say it:} \emph{porcus}
  \item \emph{Say it:} \emph{porcus}, once more
  \item \emph{Say it:} say \emph{porcus}
  \item \emph{Recall:} say the Latin for a sheep, then the Latin for a goat, then say \emph{porcus} again
\end{itemize}

\begin{tcolorbox}[breakable,colback=teal!4,colframe=teal!35!black,title={Before you move on}]
What does porcus, porcī mean? (\textquotedblleft{}A pig\textquotedblright{}.) Say it once more.
\end{tcolorbox}
